% The Chord--Tangent Law in Cases  (source §13-§16; pictorial dictionary plus
% the four structural cases. Figures reused verbatim from the source notes,
% rendered from the Sage pipeline -- including 11_correspondence.sage.)
\chapter{The Chord--Tangent Law in Cases}
\label{ch:chord_tangent_cases}

\Cref{thm:tr-main} proved once and for all that the chord--tangent operation is
the transported group law. This chapter makes the statement concrete, walking
through every structurally distinct case --- adding the identity, two distinct
points, doubling, vertical lines and inverses, and two-torsion --- and displaying
each as a picture. Every case is the single identity \(P + Q + R = O\) of
\cref{thm:lin-collinear} read in a particular configuration; nothing new is
proved, but the geometry becomes visible.

\paragraph{Colour convention.}
In the figures, rust marks the input points, plum the auxiliary line (chord,
tangent, or vertical), teal the third intersection \(R\), and gold the final
result obtained by reflecting \(R\) across the \(x\)-axis. All coordinates are
computed by the Sage scripts of \cref{app:C}; the running example is
\(E : y^2 = 4x^3 - 4x\), with \(g_2 = 4\), \(g_3 = 0\) and roots
\(e_1 = -1\), \(e_2 = 0\), \(e_3 = 1\).

\section{The Pictorial Dictionary}
\label{sec:cases-dictionary}

The isomorphism \(\Phi : \C/\Lambda \to E\) lets us watch a single point in three
synchronized pictures: its position in the fundamental parallelogram, its image
on the embedded torus, and its image on the cubic (\cref{fig:dict-overview}).
Addition on the left --- vector addition modulo \(\Lambda\) --- becomes the
chord--tangent construction on the right.

\begin{figure}[ht]
\centering
\begin{tikzpicture}[every node/.style={font=\small}]
  %=== Panel 1: Lattice Lambda ============================================
  % Scale 0.70: fundamental domain is 3.1*0.70 = 2.17 cm wide.
  % Tiles for i in {-1,0,1}, j in {0,1} (6 cells, 2 rows above x-axis).
  % Rightmost tile corner (i=1,j=1): (6.2,3.9)*0.70 = (4.34,2.73) cm in parent.
  \begin{scope}[scale=0.70]
    % Background lattice tiles
    \foreach \i in {-1,0,1}{
      \foreach \j in {0,1}{
        \pgfmathsetmacro{\xa}{\i*\latwax+\j*\latwbx}
        \pgfmathsetmacro{\ya}{\i*\latway+\j*\latwby}
        \pgfmathsetmacro{\xb}{(\i+1)*\latwax+\j*\latwbx}
        \pgfmathsetmacro{\yb}{(\i+1)*\latway+\j*\latwby}
        \pgfmathsetmacro{\xc}{(\i+1)*\latwax+(\j+1)*\latwbx}
        \pgfmathsetmacro{\yc}{(\i+1)*\latway+(\j+1)*\latwby}
        \pgfmathsetmacro{\xd}{\i*\latwax+(\j+1)*\latwbx}
        \pgfmathsetmacro{\yd}{\i*\latway+(\j+1)*\latwby}
        \draw[draw=formteal!50!black, fill=formteal!6, line width=0.35pt, opacity=0.80]
          (\xa,\ya) -- (\xb,\yb) -- (\xc,\yc) -- (\xd,\yd) -- cycle;
      }
    }
    % Lattice points
    \foreach \i in {-1,0,1,2}{
      \foreach \j in {0,1,2}{
        \pgfmathsetmacro{\xx}{\i*\latwax+\j*\latwbx}
        \pgfmathsetmacro{\yy}{\i*\latway+\j*\latwby}
        \fill[formteal!80!black] (\xx,\yy) circle (1.2pt);
      }
    }
    % Fundamental parallelogram (highlighted)
    \draw[draw=formrust!72!black, fill=formrust!13, line width=0.85pt]
      (0,0) -- (\latwax,\latway) -- (\latsxc,\latsyc) -- (\latwbx,\latwby) -- cycle;
    % Basis vectors with arrows
    \draw[-{Latex[length=1.8mm]}, formrust!78!black, line width=0.65pt]
      (0,0) -- (\latwax,\latway)
      node[midway, below, font=\scriptsize, inner sep=1.5pt] {$\omega_1$};
    \draw[-{Latex[length=1.8mm]}, formrust!78!black, line width=0.65pt]
      (0,0) -- (\latwbx,\latwby)
      node[midway, left, font=\scriptsize, inner sep=1.5pt] {$\omega_2$};
    \fill[black] (0,0) circle (1.7pt);
    \node[below left, inner sep=1.5pt, font=\scriptsize] at (0,0) {$0$};
    \node[formrust!55!black, font=\scriptsize]
      at (\latcellcx,\latcellcy) {$P_\Lambda$};
    % Caption (in scope coords, below domain)
    \node[font=\small, anchor=north] at ({(\latwax+\latwbx)/2},{\latway-0.65})
      {lattice $\Lambda\subset\C$};
  \end{scope}

  %=== Arrow: quotient map pi: C -> C/Lambda ==============================
  % Arrow placed between rightmost tile edge (~4.34 cm) and torus left edge.
  % Torus: xshift=6.30, scale=0.47 => left edge = 6.30 - 2.5*0.47 = 5.125 cm.
  \draw[-{Latex[length=2.2mm]}, formplum, line width=0.9pt] (4.55,1.0) -- (5.00,1.0);
  \node[formplum, above, font=\scriptsize, inner sep=2pt] at (4.775,1.0)
    {$\pi$};
  \node[formplum, below, font=\scriptsize, inner sep=2pt] at (4.775,1.0)
    {$/\Lambda$};

  %=== Panel 2: Torus T = C/Lambda ========================================
  % At scale 0.47: outer silhouette radius = 2.5*0.47 = 1.175 cm.
  % Scope center at parent (6.30, 1.0).
  \begin{scope}[xshift=6.30cm, yshift=1.0cm, scale=0.47]
    \drawprecisetorus
    \node[font=\small, anchor=north] at (0,-2.72)
      {torus $T=\C/\Lambda$};
  \end{scope}

  %=== Arrow: Weierstrass embedding Phi ===================================
  % Torus right edge = 6.30 + 2.5*0.47 = 7.475 cm.
  \draw[-{Latex[length=2.2mm]}, formplum, line width=0.9pt] (7.65,1.0) -- (8.15,1.0);
  \node[formplum, above, font=\scriptsize, inner sep=2pt] at (7.90,1.0)
    {$\Phi$};
  \node[formplum, below, font=\scriptsize, inner sep=2pt] at (7.90,1.0)
    {$\wp,\wp'$};

  %=== Panel 3: Elliptic curve E: y^2 = 4x^3 - 4x ========================
  % Scope center at parent (9.40, 1.0).
  % At scale 0.56: x in [-1.55, 1.69] maps to parent x in [8.532, 10.346].
  \begin{scope}[xshift=9.40cm, yshift=1.0cm, scale=0.56]
    % Axes
    \draw[->, black!50, line width=0.5pt]
      (-1.55,0) -- (2.00,0) node[right, black!55, font=\scriptsize] {$x$};
    \draw[->, black!50, line width=0.5pt]
      (0,-3.70) -- (0,3.85) node[above, black!55, font=\scriptsize] {$y$};
    % Root tick marks
    \draw[black!40, line width=0.4pt] (-1,-0.12) -- (-1,0.12);
    \draw[black!40, line width=0.4pt] ( 0,-0.12) -- ( 0,0.12);
    \draw[black!40, line width=0.4pt] ( 1,-0.12) -- ( 1,0.12);
    \node[black!60, below, font=\scriptsize] at (-1,-0.12) {$e_1$};
    \node[black!60, below right, font=\scriptsize] at (0,-0.12) {$e_2$};
    \node[black!60, below, font=\scriptsize] at (1,-0.12) {$e_3$};
    % Curve arcs from sage-computed coordinate lists
    \drawprecisecurve
    % Caption and equation label
    \node[font=\small, anchor=north] at (0.25,-3.85)
      {cubic $E\subset\CP^2$};
    \node[font=\scriptsize, formteal!50!black, anchor=north west] at (-1.50,3.85)
      {$y^2=4x^3-4x$};
  \end{scope}
\end{tikzpicture}
\caption{The dictionary \(\Lambda \leadsto T \leadsto E\). A lattice point class
\([w] = a\omega_1 + b\omega_2\) appears in the fundamental parallelogram (left),
at its embedded position on the torus (middle), and on the cubic
\(E : y^2 = 4x^3 - 4x\) (right). Computed by \texttt{sagemath/11\_correspondence.sage}
and the rendering scripts.}
\label{fig:dict-overview}
\end{figure}

\section{Two Distinct Points: the Chord}
\label{sec:cases-chord}

Let \(P, Q \in E\) be distinct finite points with \(Q \ne \ominus P\), and let
\(L\) be the line through them; it is not vertical, so it meets \(E\) in a third
finite point \(R\). Writing \(P = \Phi(u)\), \(Q = \Phi(v)\), \(R = \Phi(w)\), the
intersection divisor is \(L \cdot E = P + Q + R\), so \(\divisor(\ell/Z) =
P + Q + R - 3O\) is principal (\cref{thm:lin-principal}). Abel's law
(\cref{cor:abel-three-point}) gives \(u + v + w \equiv 0\), hence
\(w \equiv -(u+v)\) and \(R = \Phi(-(u+v))\). Since the inverse is reflection,
\[
  P \oplus Q = \Phi(u+v) = \ominus\Phi(-(u+v)) = \ominus R = -R .
\]

\begin{figure}[ht]
\centering
\begin{tikzpicture}[every node/.style={font=\small}]
  %--- Panel 1: lattice ---------------------------------------------------
  \begin{scope}[scale=0.85]
    \draw[draw=formteal!55!black, fill=formteal!8, line width=0.7pt]
      (0,0) -- (\latwax,\latway) -- (\latsxc,\latsyc) -- (\latwbx,\latwby) -- cycle;
    \draw[-{Latex[length=1.5mm]}, formteal!60!black, line width=0.5pt, opacity=0.75]
      (0,0) -- (\latwax,\latway) node[midway, below, font=\scriptsize] {$\omega_1$};
    \draw[-{Latex[length=1.5mm]}, formteal!60!black, line width=0.5pt, opacity=0.75]
      (0,0) -- (\latwbx,\latwby) node[midway, left, font=\scriptsize] {$\omega_2$};
    % Origin
    \fill[black] (\corOriginLx,\corOriginLy) circle (1.6pt)
      node[below left, inner sep=1pt, font=\scriptsize] {$0$};
    % u, v, sum
    \fill[formrust!85!black] (\corULx,\corULy) circle (1.9pt);
    \node[formrust!60!black, below right, inner sep=1pt, font=\scriptsize] at (\corULx,\corULy) {$u$};
    \fill[formrust!85!black] (\corVLx,\corVLy) circle (1.9pt);
    \node[formrust!60!black, above right, inner sep=1pt, font=\scriptsize] at (\corVLx,\corVLy) {$v$};
    \fill[formgold!85!black] (\corSumLiftLx,\corSumLiftLy) circle (1.9pt);
    \node[formgold!50!black, above right, inner sep=1pt, font=\scriptsize]
      at (\corSumLiftLx,\corSumLiftLy) {$u{+}v$};
    % Translation arrows from origin -> u -> u+v (forming a "+v" shift)
    \draw[-{Latex[length=1.8mm]}, draw=formgold!70!black, line width=0.55pt, dashed]
      (\corULx,\corULy) -- (\corSumLiftLx,\corSumLiftLy);
    \draw[-{Latex[length=1.8mm]}, draw=formgold!70!black, line width=0.55pt, dashed]
      (\corVLx,\corVLy) -- (\corSumLiftLx,\corSumLiftLy);
    \node[font=\small] at (1.55,-0.65) {lattice};
  \end{scope}
  %--- Arrow lattice -> torus ---------------------------------------------
  \draw[-{Latex[length=2.4mm]}, formplum, line width=0.9pt] (3.05,1.0) -- (3.95,1.0);
  \node[formplum] at (3.5,1.35) {$\pi$};
  %--- Panel 2: torus -----------------------------------------------------
  \begin{scope}[xshift=5.65cm, yshift=1.1cm, scale=0.45]
    \drawprecisetorus
    \fill[black] (\corOriginTx,\corOriginTy) circle (2.6pt);
    \node[anchor=south west, inner sep=1pt, font=\scriptsize] at (\corOriginTx,\corOriginTy) {$[0]$};
    \fill[formrust!85!black] (\corUTx,\corUTy) circle (3.6pt);
    \node[formrust!60!black, anchor=south, inner sep=1pt, font=\scriptsize]
      at ({\corUTx},{\corUTy+0.15}) {$[u]$};
    \fill[formrust!85!black] (\corVTx,\corVTy) circle (3.6pt);
    \node[formrust!60!black, anchor=south east, inner sep=1pt, font=\scriptsize]
      at ({\corVTx-0.05},{\corVTy+0.05}) {$[v]$};
    \fill[formgold!85!black] (\corSumRedTx,\corSumRedTy) circle (3.6pt);
    \node[formgold!50!black, anchor=east, inner sep=1pt, font=\scriptsize]
      at ({\corSumRedTx-0.15},{\corSumRedTy}) {$[u{+}v]$};
    \node[font=\small] at (0,-2.85) {torus $T=\C/\Lambda$};
  \end{scope}
  %--- Arrow torus -> curve -----------------------------------------------
  \draw[-{Latex[length=2.4mm]}, formplum, line width=0.9pt] (7.4,1.0) -- (8.3,1.0);
  \node[formplum] at (7.85,1.35) {$\Phi$};
  %--- Panel 3: curve -----------------------------------------------------
  \begin{scope}[xshift=10.0cm, yshift=-0.95cm, scale=0.55]
    \draw[->, black!50, line width=0.5pt] (-1.6,0) -- (2.9,0) node[right, black!55, font=\scriptsize] {$x$};
    \draw[->, black!50, line width=0.5pt] (0,-3.4) -- (0,3.6) node[above, black!55, font=\scriptsize] {$y$};
    % Curve arcs and branch points from sagemath/12_curve_sample.sage
    \drawprecisecurve
    \draw[draw=formplum, line width=0.85pt]
      (\corCurveCLeftx,\corCurveCLefty) -- (\corCurveCRightx,\corCurveCRighty);
    \draw[dashed, line width=0.55pt, formgold!75!black]
      (\corCurveCRx,\corCurveCRy) -- (\corCurveCSumx,\corCurveCSumy);
    \fill[formrust!85!black] (\corCurveCPx,\corCurveCPy) circle (2.3pt);
    \node[formrust!60!black, above left, font=\scriptsize] at (\corCurveCPx,\corCurveCPy) {$P$};
    \fill[formrust!85!black] (\corCurveCQx,\corCurveCQy) circle (2.3pt);
    \node[formrust!60!black, above right, font=\scriptsize] at (\corCurveCQx,\corCurveCQy) {$Q$};
    \fill[formteal!85!black] (\corCurveCRx,\corCurveCRy) circle (2.2pt);
    \node[formteal!60!black, above left, font=\scriptsize] at (\corCurveCRx,\corCurveCRy) {$R$};
    \fill[formgold!85!black] (\corCurveCSumx,\corCurveCSumy) circle (2.6pt);
    \node[formgold!50!black, below left, font=\scriptsize] at (\corCurveCSumx,\corCurveCSumy) {$P{+}Q$};
    \node[font=\small] at (0.65,-3.9) {curve $E$};
  \end{scope}
\end{tikzpicture}
\caption{Adding two distinct points. On the lattice (left) the gold arrows show
\(u \to u+v \leftarrow v\); on the torus (middle) both land at \([u+v]\); on the
cubic (right) the chord \(PQ\) meets \(E\) at \(R\), and \(P \oplus Q = -R\) is
its reflection. Coordinates from \texttt{sagemath/11\_correspondence.sage}.}
\label{fig:cases-chord}
\end{figure}

\begin{keyidea}{Why the sign appears}
The line gives \(P + Q + R = O\), not \(P + Q = R\). Solving forces
\(P \oplus Q = -R\): the final reflection is exactly this minus sign, traced back
to the torus inverse \(z \mapsto -z\).
\end{keyidea}

\section{Doubling: the Tangent}
\label{sec:cases-double}

When \(P = Q\) the chord degenerates to the tangent line at \(P\), the limit of
the chord as \(Q \to P\). The tangent meets \(E\) at \(P\) with intersection
multiplicity two (\cref{ex:lin-tangent}), so \(L \cdot E = 2P + R\). With
\(P = \Phi(u)\), \(R = \Phi(w)\), the principal divisor \(2P + R - 3O\) and
Abel's law give \(2u + w \equiv 0\), so \(w \equiv -2u\), \(R = \Phi(-2u)\), and
\[
  2P = \Phi(2u) = \ominus R = -R .
\]

\begin{figure}[ht]
\centering
\begin{tikzpicture}[scale=1.1, every node/.style={font=\small}]
  \draw[->, black!50, line width=0.5pt] (-1.8,0) -- (3.0,0) node[right, black!55] {$x$};
  \draw[->, black!50, line width=0.5pt] (0,-3.4) -- (0,3.5) node[above, black!55] {$y$};
  % Sage curve y^2 = 4x^3 - 4x
  % Curve arcs and branch points from sagemath/12_curve_sample.sage
  \drawprecisecurve
  \foreach \rt in {\rooteone,\rootetwo,\rootethree}
    \fill[formteal!60!black] (\rt,0) circle (1.3pt);
  % Sage-computed tangent line at P
  \draw[draw=formplum, line width=0.85pt]
    (\dbllefta,\dblleftb) -- (\dblrighta,\dblrightb);
  % Reflection line
  \draw[dashed, line width=0.45pt, formgold!75!black]
    (\dblRx,\dblRy) -- (\dbltwoPx,\dbltwoPy);
  % P, R, 2P
  \fill[formrust!85!black] (\dblPx,\dblPy) circle (2.1pt);
  \node[formrust!60!black, above right] at (\dblPx,\dblPy) {$P$};
  \draw[draw=formrust!75!black, line width=0.6pt] (\dblPx,\dblPy) circle (4.5pt);
  \fill[formteal!85!black] (\dblRx,\dblRy) circle (2pt);
  \node[formteal!60!black, below right] at (\dblRx,\dblRy) {$R$};
  \fill[formgold!85!black] (\dbltwoPx,\dbltwoPy) circle (2.4pt);
  \node[formgold!50!black, above right] at (\dbltwoPx,\dbltwoPy) {$2P=-R$};
  % Annotation
  \node[draw=black!40, rounded corners=2pt, fill=formivory, align=left, anchor=west,
        inner sep=3pt, font=\scriptsize] at (-1.85,-2.4)
    {Sage: $P=(\dblPx,\dblPy)$\\
     tangent slope $m=\dblslope$\\
     Vieta with double root $\Rightarrow$\\
     $x_R=\tfrac{m^2}{4}-2x_P$\\
     $\Rightarrow R=(\dblRx,\dblRy)$};
\end{tikzpicture}
\caption{Doubling a point. The tangent to \(E\) at \(P\) (plum) meets \(E\) again
at \(R\), and \(2P = -R\) is its reflection (gold). The slope is
\(m = (12x_P^2 - g_2)/(2y_P)\); see \cref{ch:coordinate_formulas}. Coordinates
from \texttt{sagemath/05\_tangent\_double.sage}.}
\label{fig:cases-tangent}
\end{figure}

\section{The Identity and Inverses: Vertical Lines}
\label{sec:cases-vertical}

The identity is \(O = \Phi(0)\). A line through a finite point \(P = (x,y)\) and
\(O\) is the vertical line \(X = xZ\), which also passes through \(-P = (x,-y)\).
Its intersection divisor is \(P + (\ominus P) + O\), so \cref{thm:lin-collinear}
gives \(P \oplus (\ominus P) \oplus O = O\). Reading this two ways: with the third
point \(\ominus P\) it says \(P \oplus O = P\) (the identity law); directly it says
\(P \oplus (\ominus P) = O\) (inverses), confirming that \(\ominus P = (x,-y)\) is
the group inverse.

\begin{figure}[ht]
\centering
\begin{tikzpicture}[scale=1.1, every node/.style={font=\small}]
  \draw[->, black!50, line width=0.5pt] (-1.8,0) -- (3.0,0) node[right, black!55] {$x$};
  \draw[->, black!50, line width=0.5pt] (0,-3.4) -- (0,3.5) node[above, black!55] {$y$};
  % Curve arcs and branch points from sagemath/12_curve_sample.sage
  \drawprecisecurve
  \foreach \rt in {\rooteone,\rootetwo,\rootethree}
    \fill[formteal!60!black] (\rt,0) circle (1.3pt);
  % Sage-computed vertical line and the two intersections
  \draw[draw=formplum, line width=0.85pt] (\vertx,-3.3) -- (\vertx,3.3);
  \fill[formrust!85!black] (\vertx,\verty) circle (2pt);
  \node[formrust!60!black, right] at (\vertx,\verty) {$P=(x_0,y_0)$};
  \fill[formgold!85!black] (\vertx,{-\verty}) circle (2pt);
  \node[formgold!50!black, right] at (\vertx,{-\verty}) {$-P=(x_0,-y_0)$};
  \draw[draw=formplum, dashed, line width=0.55pt, -{Latex[length=2.0mm]}]
    (\vertx,3.3) to[bend right=14] ({\vertx+0.55},3.5);
  \node[formplum, font=\scriptsize, right] at ({\vertx+0.55},3.5) {$\to O=[0{:}1{:}0]$};
  \node[draw=black!40, rounded corners=2pt, fill=formivory, align=left, anchor=west,
        inner sep=3pt, font=\scriptsize] at (-1.85,-2.5)
    {Sage: $x_0=\vertx,\;y_0=\verty$\\
     Vertical line $x=x_0$ meets $E$\\
     at $P$, $-P$ and $O$ at infinity,\\
     so $P+(-P)+O=O$.};
\end{tikzpicture}
\caption{Vertical lines give inverses. The line \(x = x_0\) meets \(E\) at
\(P = (x_0, y_0)\), \(-P = (x_0, -y_0)\), and \(O\) at infinity, so
\(P \oplus (\ominus P) = O\). Coordinates from \texttt{sagemath/06\_two\_torsion.sage}.}
\label{fig:cases-vertical}
\end{figure}

\section{Two-Torsion Points}
\label{sec:cases-torsion}

A point is its own inverse, \(P = \ominus P\), exactly when \(2P = O\): the
\emph{two-torsion}. In affine coordinates \(P = \ominus P\) means \((x,y) =
(x,-y)\), so \(y = 0\); the nonzero two-torsion points are therefore the three
\(x\)-axis crossings
\[
  4x^3 - g_2 x - g_3 = 0, \qquad\text{i.e.}\qquad x \in \{e_1, e_2, e_3\},
\]
the roots of \cref{prop:wp-roots}. At each the tangent is vertical, so doubling
gives \(2P = O\) directly. Together with \(O\) they form
\[
  E[2] = \{O, T_1, T_2, T_3\} \cong (\Z/2\Z)^2,
\]
the images under \(\Phi\) of the two-torsion of \cref{ex:tg-two-torsion}.

\begin{figure}[ht]
\centering
\begin{tikzpicture}[scale=1.1, every node/.style={font=\small}]
  \draw[->, black!50, line width=0.5pt] (-1.8,0) -- (3.0,0) node[right, black!55] {$x$};
  \draw[->, black!50, line width=0.5pt] (0,-3.4) -- (0,3.5) node[above, black!55] {$y$};
  % Curve arcs and branch points from sagemath/12_curve_sample.sage
  \drawprecisecurve
  % Sage-computed real 2-torsion x-coords
  \fill[formrust!85!black] (\torstwoone,0) circle (2.4pt);
  \node[formrust!60!black, above] at (\torstwoone,0.10) {$T_1$};
  \fill[formrust!85!black] (\torstwotwo,0) circle (2.4pt);
  \node[formrust!60!black, above] at (\torstwotwo,0.10) {$T_2$};
  \fill[formrust!85!black] (\torstwothree,0) circle (2.4pt);
  \node[formrust!60!black, above] at (\torstwothree,0.10) {$T_3$};
  % Vertical tangents at the roots
  \draw[draw=formplum, dashed, line width=0.5pt] (\torstwoone,-3.3) -- (\torstwoone,3.3);
  \draw[draw=formplum, dashed, line width=0.5pt] (\torstwotwo,-3.3) -- (\torstwotwo,3.3);
  \draw[draw=formplum, dashed, line width=0.5pt] (\torstwothree,-3.3) -- (\torstwothree,3.3);
  \node[formgold!50!black, font=\small] at (2.5,3.1) {$O=[0{:}1{:}0]$};
  \node[draw=black!40, rounded corners=2pt, fill=formivory, align=left, anchor=west,
        inner sep=3pt, font=\scriptsize] at (-1.85,-2.6)
    {Sage roots of $4x^3-4x=0$:\\
     $T_1=(\torstwoone,0)$\\
     $T_2=(\torstwotwo,0)$\\
     $T_3=(\torstwothree,0)$\\
     $E[2]=\{O,T_1,T_2,T_3\}\cong(\Z/2)^2$};
\end{tikzpicture}
\caption{The two-torsion. The nonzero two-torsion points \(T_1, T_2, T_3\) are
the \(x\)-axis crossings \((e_i, 0)\), each with a vertical tangent; with \(O\)
they form \(E[2] \cong (\Z/2\Z)^2\). Roots of \(4x^3 - 4x\) from
\texttt{sagemath/06\_two\_torsion.sage}.}
\label{fig:cases-torsion}
\end{figure}

The three nonzero two-torsion points are collinear --- they all lie on the
\(x\)-axis \(y = 0\), itself a line meeting \(E\) at \((e_1,0), (e_2,0), (e_3,0)\).
By \cref{thm:lin-collinear}, \(T_1 \oplus T_2 \oplus T_3 = O\), which is the
statement that the three nonzero elements of \((\Z/2\Z)^2\) sum to zero.

\section{Summary}
\label{sec:cases-summary}

Each case of the chord--tangent law is one reading of \(P + Q + R = O\). For two
distinct points the line is a chord and \(P \oplus Q = -R\)
(\cref{sec:cases-chord}); for a repeated point it is the tangent and \(2P = -R\)
(\cref{sec:cases-double}); a vertical line through \(P\) and \(O\) yields the
identity law and exhibits \(\ominus P = (x,-y)\) (\cref{sec:cases-vertical}); and
the points fixed by reflection are the two-torsion \(E[2] \cong (\Z/2\Z)^2\)
(\cref{sec:cases-torsion}). The pictures make the dictionary
\(\Lambda \leadsto T \leadsto E\) literal. \Cref{ch:coordinate_formulas} converts
these geometric rules into closed-form coordinate formulas.

\section*{Exercises}
\addcontentsline{toc}{section}{Exercises}

\begin{exercise}
\label{ex:cases-chord-vieta}
Let \(P = (x_P, y_P)\), \(Q = (x_Q, y_Q)\) be distinct with chord \(y = mx + k\).
Substitute into \(y^2 = 4x^3 - g_2 x - g_3\) and use Vieta's formulas on the
resulting cubic in \(x\) to show the third intersection has
\(x_R = \tfrac{m^2}{4} - x_P - x_Q\). Then \(P \oplus Q = (x_R, -(m x_R + k))\).
\end{exercise}

\begin{exercise}
\label{ex:cases-tangent-slope}
For doubling, differentiate \(y^2 = 4x^3 - g_2 x - g_3\) implicitly to get the
tangent slope \(m = (12 x_P^2 - g_2)/(2 y_P)\), and show the doubled point has
\(x_{2P} = \tfrac{m^2}{4} - 2 x_P\). What happens to \(m\) when \(y_P = 0\), and
how does that match \(2P = O\)?
\end{exercise}

\begin{exercise}
\label{ex:cases-two-torsion-sum}
Show that the \(x\)-axis is the line through the three nonzero two-torsion points,
and deduce \(T_1 \oplus T_2 \oplus T_3 = O\) from \cref{thm:lin-collinear}.
Verify this is consistent with \(E[2] \cong (\Z/2\Z)^2\), where the three nonzero
elements sum to zero.
\end{exercise}

\begin{exercise}
\label{ex:cases-flex-three-torsion}
A point \(P\) is \emph{three-torsion}, \(3P = O\), if and only if \(P\) is an
inflection point of \(E\) (the tangent at \(P\) meets \(E\) with multiplicity
three). Using \(L \cdot E = 3P\) and \cref{thm:lin-collinear}, prove this, and note
\(O\) is one such flex (\cref{prop:emb-flex}).
\end{exercise}
